\subsection{Rayleigh--Taylor 不安定性：線形成長から非線形形状へ}
\label{sec:val_rayleigh_taylor}
% ═══════════════════════════════════════════════════════════════════

\subsubsection{設定と理論基準}

Rayleigh--Taylor ベンチマークでは，水（重相，上）／シリコーン油 5cSt（軽相，下）の
二層配置に小振幅摂動を与え，重力下で不安定成長させる．
本稿の記号規約では，
液相記号を下層シリコーン油，
気相記号を上層水に対応させる．
これは物理相の密度順序を優先するための記号規約であり，
pressure-jump の向き $j_{gl}=p_g-p_l$ は他ケースと同じままである．

線形理論では，小振幅擾乱 $A_0$ は
\begin{equation}
  A(t) = A_0 \, e^{\gamma t},
  \qquad
  \gamma = \sqrt{\mathrm{At}\, g\, k},
  \qquad
  \mathrm{At} = \frac{\rho_{\text{w}}-\rho_{\text{oil}}}{\rho_{\text{w}}+\rho_{\text{oil}}},
  \qquad
  k = \frac{2\pi m}{L_x}
  \label{eq:rt_linear_growth}
\end{equation}
を第一基準に置く~\cite{Chandrasekhar1961}．
本設定では
$\mathrm{At}\approx0.0256$，$k=4\pi$，
$\gamma\approx1.78\,\mathrm{s^{-1}}$ であり，
$T=7.0$ は線形領域を越えて非線形 mushroom 形成を観察する時間である．

\begin{itemize}[noitemsep,leftmargin=2em]
  \item グリッド：$[0,0.5]\times[0,2.0]$，$128\times512$，
        壁面境界，interface-fitted $\alpha=2$
  \item 物性：下層シリコーン油 $\rho=950,\mu=5\times10^{-3}$；
        上層水 $\rho=1000,\mu=10^{-3}$
  \item 表面張力・重力：$\sigma=0.040$，$g=9.81$
  \item 初期界面：$y_0=1.0$，$A_0=0.01$，$m=1$，$L_x=0.5$
  \item 時間：$T=7.0$，理論 CFL multiplier $1.0$
  \item 出力：\texttt{interface\_amplitude}，
        \texttt{kinetic\_energy}，\texttt{volume\_conservation}，
        $\psi$/速度/圧力 snapshot
\end{itemize}

\subsubsection{結果と判定}

早期窓では，\texttt{interface\_amplitude} が
式~\eqref{eq:rt_linear_growth} の対数線形増幅と整合する．
その後，snapshot 上では spike/bubble 構造が発達し，
運動エネルギーは線形増幅に伴って増加したのち，非線形領域で有界な推移へ移る．
体積保存が保たれることで，mushroom 形状は相体積の数値漏れではなく，
浮力不安定性の形状発展として読める．

本ケースの判定は，
「pressure-jump と hydrostatic-balanced buoyancy を同じ stack で使ったまま，
安定振動だけでなく不安定成長も正しい向きに発生させられる」
である．
毛管波・振動液滴が表面張力主導の復元問題であるのに対し，
Rayleigh--Taylor は浮力主導の成長問題であり，
同じ面速度空間の契約が両方を支えることを示した．

% ═══════════════════════════════════════════════════════════════════
